\chapter{真空のエネルギー、インフレーション、原始揺らぎ}

\section{真空のエネルギーとド・ジッター空間}

平坦な空間ではフリードマン方程式は
\begin{equation}
\left(\frac{\dot A}{A}\right)^2=\frac{8\pi G}{3}\rho.
\end{equation}
である。物質については、\(\rho\)は\(A^{-3}\)のように減少する。放射については、\(\rho\)は\(A^{-4}\)のように減少する。真空のエネルギーについては、それは薄まらない：
\begin{equation}
\rho=\rho_0=\text{constant}.
\end{equation}
宇宙がこの真空のエネルギーのみを含むならば、右辺は定数である。その定数を\(H^2\)と書くと、
\begin{equation}
H^2=\frac{8\pi G}{3}\rho_0,
\qquad
H=\sqrt{\frac{8\pi G}{3}\rho_0}.
\end{equation}
となる。するとハッブル定数は真に定数である：
\begin{equation}
\frac{\dot A}{A}=H.
\end{equation}
したがって
\begin{equation}
\dot A=HA,
\end{equation}
であり、解は指数関数的である：
\begin{equation}
A(t)=C e^{Ht}.
\end{equation}
束縛されていない遠方の銀河の間の距離は指数関数的に増大する。これがド・ジッター空間である。

これらの座標では幾何学は
\begin{equation}
d\tau^2=dt^2-e^{2Ht}\left(dx^2+dy^2+dz^2\right).
\end{equation}
である。ここでの真空のエネルギーは正である。平坦な空間ではそうならざるを得ない。なぜならフリードマン方程式の左辺は非負だからである。負の真空のエネルギーは曲率の項を含めるときにのみ現れうるが、ここで追っているのはそのような場合ではない。

アインシュタインが宇宙定数を導入したのは別の目的のためであった。すなわち、通常の物質と釣り合いをとり、宇宙を静的 に保つためであった。

\section{現在の加速する宇宙}

ド・ジッター空間は現在の宇宙の晩期における傾向でもある。以前はスケール因子は物質優勢の振る舞い
\begin{equation}
A(t)\propto t^{2/3}.
\end{equation}
に従っていた。数十億年前に曲線は上向きに曲がり、より急速に増大し始めた。他の形態のエネルギーが依然として存在するため、まだ完全な指数関数ではないが、傾向は指数関数的なものへと向かっている。

ここで加速とは
\begin{equation}
\ddot A>0.
\end{equation}
を意味する。これは晩期の膨張の観測された特徴である。

次に地平線を考えよう。ハッブルの法則は
\begin{equation}
v=HD.
\end{equation}
と言う。後退速度が光速になる距離では、
\begin{equation}
D=\frac{c}{H}.
\end{equation}
である。等価的に、これを\(HD=c\)と書き直すことができる。\editorialnote{これは発話された関係\(D=c/H\)を編集上代数的に言い換えたものである。図~\ref{fig:lecture07-horizon}の右側の黒板の上の行は一部しか判読できない。}

\begin{figure}[ht]
\centering
\includegraphics[width=0.62\textwidth]{lecture_07_figure_02.png}
\caption{ド・ジッター膨張における宇宙論的地平線の関係\(D=c/H\)。}
\label{fig:lecture07-horizon}
\lectureframe{00:30:37}
\end{figure}

この距離は百億から百五十億光年の程度であり、後者を\(c\)で距離に換算すれば、宇宙の年齢に数値的に近い。

\begin{classroomqa}
\lecturetimestamp{00:14:10}
\audiencequestion{地平線のスケールと宇宙の年齢がほぼ等しいのは、接線がたまたま原点を通る場合に生じることではないのか。}
\lecturerresponse{それは\(A(t)\propto t\)に対しては正確に成り立つであろう。一般には成り立たない。ド・ジッター空間では\(H\)は固定されたままで年齢は増え続けるので、現在の近似的な等しさは数学的恒等式ではない。}
\end{classroomqa}

\section{真空のエネルギーの意味}

真空のエネルギーはここでは通常の量子力学の言葉で記述される。調和振動子についてエネルギーは模式的には
\begin{equation}
E\sim p^2+x^2.
\end{equation}
である。古典的には基底状態は平衡点であり、
\begin{equation}
x=0,\qquad p=0,\qquad E=0.
\end{equation}
である。量子力学は\(x\)と\(p\)の両方が正確にゼロに固定されることを許さない。不確定性原理が揺らぎを強制するので、基底状態のエネルギーは正である。振動子については
\begin{equation}
E_0=\frac{1}{2}\hbar\omega.
\end{equation}
が得られる。

同じ論点は次に場に持ち越される。宇宙は調和振動子で満ちている。電磁場の振動、そして同様に他のあらゆる場の振動である。各々の場は零点エネルギーを持つ。電磁場についてエネルギー密度は模式的には
\begin{equation}
u_{\mathrm{EM}}\sim E^2+B^2.
\end{equation}
である。実光子が存在しなくても、不可避の零点振動が残る。これが真空のエネルギーの一つの説明である。ここで必要な唯一の性質は、このエネルギーが存在し、宇宙が膨張しても薄まらないということである。

\begin{classroomqa}
\lecturetimestamp{00:16:15}
\audiencequestion{この真空のエネルギーとは何であり、物質と放射を取り除いてもなぜ何かが残るのか。}
\lecturerresponse{真空は量子場の基底状態であり、量子基底状態はゼロのエネルギーを持つ必要はないからである。調和振動子は正の基底状態エネルギーを持ち、場は振動子の集まりとして振る舞う。したがって実粒子が存在しなくても零点エネルギーが残る。}
\end{classroomqa}

これは三度の背景放射ではない。マイクロ波背景放射は通常の放射である。そのエネルギー密度は光子が薄まり、各光子が赤方偏移するために減少する。真空のエネルギーは、実光子が取り除かれた後に残る残留零点エネルギーである。

\section{\(H^{-1}\)の意味}

ここで、逆ハッブル定数と宇宙の年齢との関係に戻ろう。もし
\begin{equation}
A(t)\propto t,
\end{equation}
ならば
\begin{equation}
H=\frac{\dot A}{A}=\frac{1}{t}.
\end{equation}
である。より一般に、もし
\begin{equation}
A(t)\propto t^p,
\end{equation}
ならば
\begin{equation}
H=\frac{p}{t}.
\end{equation}
である。物質優勢の場合 \(p=2/3\) では、
\begin{equation}
H=\frac{2}{3t}.
\end{equation}
となる。したがって、冪乗則宇宙論では \(H^{-1}\) は年齢と同じ程度の大きさであるが、正確に年齢に等しいわけではない。

ド・ジッター空間は異なる。そこでは
\begin{equation}
A(t)=C e^{\gamma t},
\end{equation}
であり、したがって
\begin{equation}
H=\frac{\dot A}{A}=\gamma,
\end{equation}
は真の定数である。ここで \(H^{-1}\) は宇宙の年齢ではない。それは指数関数的成長の時定数、すなわちスケール因子が \(e\) 倍されるのに要する時間である。口頭の議論では繰り返し ``倍'' と言った後で訂正しているが、正確な記述は \(e\) 倍である。今日、その時定数は百億年の程度である。

\begin{classroomqa}
\lecturetimestamp{00:23:46}
\audiencequestion{宇宙の年齢と逆ハッブル定数についての点をもう一度説明していただけますか。}
\lecturerresponse{冪乗則では \(H=p/t\) が得られるので、\(H^{-1}\) は1の程度の因子を除いて年齢を追跡する。しかしド・ジッター空間では \(H\) は一定である一方、年齢は増え続けるので、\(H^{-1}\) は指数関数的な時定数にすぎず、年齢ではない。}
\end{classroomqa}

\section{エネルギーと宇宙の地平線}

真空エネルギーの密度が宇宙の膨張中も一定に保たれるならば、真空エネルギーの総量が増加しているのかどうかを問うのは自然である。ここで与えられる答えには限界がある。一般相対性理論ではエネルギーの収支に重力の寄与が含まれるため、全エネルギーという概念は初等的な力学におけるよりも微妙である。実用的な要点は地平線に関わるものである。

厳密なド・ジッター空間では、宇宙の地平線は固定された大きさをもつ。
\begin{equation}
D=\frac{c}{H}
=\frac{c}{\sqrt{8\pi G\rho_0/3}}.
\end{equation}
したがって、その地平線の内側の体積は固定されている。その内部の真空エネルギーは
\begin{equation}
E_{\mathrm{hor}}=\frac{4\pi}{3}D^3\rho_0.
\end{equation}
\editorialnote{これは 00:31:06--00:32:23 における口頭の計算を要約したものである。すなわち、地平線の大きさをとり、球の体積を求め、\(\rho_0\) を掛ける。}
\(D\) と \(\rho_0\) はどちらも定数であるから、この量は時間とともに変化しない。ある観測者がいずれ見ることのできる宇宙の部分について、厳密なド・ジッター空間ではアクセス可能な真空エネルギーは固定されたままである。

物事は、ブラックホールの地平線を通り抜けるのとほぼ同じ意味で、宇宙の地平線を通り抜けることもできる。ひとたび外に出れば、それらは我々との因果的接触から外れる。

\begin{classroomqa}
\lecturetimestamp{00:29:27}
\audiencequestion{真空エネルギーの密度が宇宙の膨張中も一定に保たれるならば、余分なエネルギーはどこから来るのですか。}
\lecturerresponse{講義はニュートン的な箱の言葉ではなく、一般相対性理論の言葉で答える。厳密なド・ジッター空間では地平線の大きさは固定されたままであるから、いずれ見ることのできる領域の内側の真空エネルギーもまた固定されたままである。}
\end{classroomqa}

\section{なぜインフレーションが導入されたか}

次に、同じド・ジッターの機構がはるか過去へと押しやられるが、はるかに短い時定数、程度として
\begin{equation}
H^{-1}\sim 10^{-32}\,\mathrm{s}.
\end{equation}
をもつ。これがインフレーションの考え方である。すなわち、直接観測可能なものより前に、宇宙は指数関数的膨張の期間を経験した。

基本的な動機は宇宙の平坦性と滑らかさである。空間的断面は非常に平坦に見える。大きなスケールでは宇宙はまた非常に一様である。さらに顕著なことに、最終散乱面はすでに滑らかであったので、平滑化は後の可視的な歴史の間にのみ起こったはずがない。

しぼんだ風船はしわだらけである。十分に膨らませれば、見える任意の部分は滑らかになる。レーズンやプルーンの類推も同じ点を示す。しわのある物体は、巨大な因子で引き伸ばされると、見える任意の部分で滑らかに見える。インフレーションはまさにそれを行ったはずである。

これはここで用いられている形の地平線問題にも答える。最終散乱面の広く離れた領域はどちらも我々に見えるが、光が一方から他方へ移動する時間は後の可視的な歴史の間にはなかった。したがって、それらのほぼ一様な密度と温度は、より小さな一つの領域が巨大に引き伸ばされた前史を示唆する。

インフレーションには混ざった美点がある。それは初期条件への感受性を消し去り、同じ理由で宇宙がどのように始まったかについての情報の多くを消し去る。

平坦性の問題はここでより鋭い形で再び現れる。観測的には曲率パラメータ \(k\) は、約1パーセントかそれ以上の水準でゼロのように見える。幾何学は可視領域にわたってユークリッド的に見える。インフレーションは、曲率半径を我々が見ることのできる領域よりはるかに大きくすることによってこれを説明する。

\begin{classroomqa}
\lecturetimestamp{00:40:05}
\audiencequestion{我々はインフレーションの前に平衡を仮定しているのか。}
\lecturerresponse{いいえ。インフレーションの要点はほとんど逆である。それ以前にあった構造はどんなものであれ、あまりにも引き伸ばされるため、可視領域は以前の詳細の記憶を失う。}
\end{classroomqa}

\section{スカラー場とそのポテンシャル}

有効真空エネルギーが初期宇宙において現在よりもはるかに大きかった可能性があることを説明するために、スカラー場を導入する。ベクトルである電場や磁場とは異なり、スカラー場には方向がない。それは単に各点における値を持つ：
\begin{equation}
\phi=\phi(x,y,z,t).
\end{equation}
この場は仮説的なものである。自然界の既知のいかなる場とも同一視されず、その量子も既知のいかなる粒子とも同一視されない。

ここで重要となるそのエネルギーの部分はポテンシャルエネルギーであり、
\begin{equation}
V(\phi).
\end{equation}
``ポテンシャルエネルギー''とは、場の時間微分ではなく、場の値そのものに付随するエネルギーを意味する。\(\phi\)の異なる値は異なるエネルギーを持ちうる。

\begin{figure}[ht]
\centering
\includegraphics[width=0.68\textwidth]{lecture_07_figure_03.png}
\caption{定性的なスカラーポテンシャル\(V(\phi)\)。異なる場の値が異なるポテンシャルエネルギーを担うことを示す。}
\label{fig:lecture07-generic-potential}
\lectureframe{00:47:08}
\end{figure}

\begin{classroomqa}
\lecturetimestamp{00:46:34}
\audiencequestion{\(V(\phi)\)の図における横軸は何か。}
\lecturerresponse{横軸は場の値\(\phi\)であり、縦軸はポテンシャルエネルギー\(V\)である。}
\end{classroomqa}

場が空間的にも時間的にも一定であるか、あるいは関心のある領域にわたってほぼ一定であるならば、\(V(\phi)\)は真空エネルギーとまったく同じように振る舞う。それは希薄化しない。したがって、以前のド・ジッターの式は\(\rho_0\)を\(V(\phi)\)で置き換えて再利用される。特に、
\begin{equation}
H=\sqrt{\frac{8\pi G}{3}V(\phi)}.
\end{equation}

\(\phi\)が場所ごとに変化するならば、異なる領域は異なる時定数でインフレーションを起こす。しかし引き伸ばしが十分に大きければ、可視のパッチはどれもほぼ一様になる。

\begin{classroomqa}
\lecturetimestamp{01:05:11}
\audiencequestion{ここで\(\phi\)は時間パラメータのように振る舞っているのか。}
\lecturerresponse{この議論では、場がポテンシャルを転がり落ちる間は時間の単調関数であるという限られた意味において、そう振る舞う。}
\end{classroomqa}

\section{プラトー、崖、再加熱}

この物語に必要なポテンシャルは非常に平坦になるが、完全に平坦ではなく、その後崖の縁に達し、急激に下降し、完全にはゼロにならず、再び上昇する。

\begin{figure}[ht]
\centering
\includegraphics[width=0.68\textwidth]{lecture_07_figure_04.png}
\caption{高いプラトー、急な下降、より低い極小を持つインフレーション・ポテンシャル。}
\label{fig:lecture07-inflationary-potential}
\lectureframe{01:02:30}
\end{figure}

この形状はきわめて人為的であるが、宇宙論的な物語がそのようなものを必要としているように見えるため導入される。

場はプラトーの高いところから始まる。傾きが小さく、膨張に伴う有効粘性が大きいため、\(\phi\)はごくゆっくりとしか変化しない。テーブル上を滑る物体よりも、蜂蜜の中をゆっくりと落下する岩の方が意図された描像に近い。宇宙が速く膨張するほど、この粘性は大きくなる。

\begin{classroomqa}
\lecturetimestamp{00:59:07}
\audiencequestion{何が\(\phi\)を動かしているのか。}
\lecturerresponse{ポテンシャルの傾きがそれを押している。傾きは小さく、宇宙論的粘性は大きいので、ドリフトは速くではなく遅い。}
\end{classroomqa}

\(\phi\)がプラトーに沿って少しずつ進む間、\(V(\phi)\)はごくわずかしか変化しない。そのため宇宙は真空エネルギーが一定であるかのようにほとんどインフレーションを起こす。この段階の時定数は程度として
\begin{equation}
H^{-1}\sim 10^{-32}\,\mathrm{s}.
\end{equation}
大まかな口頭の見積もりでは、プラトーの高さはプランク単位で\(10^{-14}\)付近である。\editorialnote{これは01:14:38--01:14:47における口頭の``最有力''見積もりであり、講義で導出された値ではない。}

インフレーションはどれくらい続いたのか。それは知られていない。プラトーの幅と浅さが知られていないからである。知られているのは、宇宙が少なくともおよそ五十回\(e\)倍されたということであり、通常引用される数は六十である。\editorialnote{01:14:05--01:19:27における口頭のe-foldingの議論の編集上の言い換え。} これらはe-foldingである：各e-foldは\(e\)倍する一回の乗算である。

場が縁に達すると、ポテンシャルエネルギーは急速に変化する。場は落下し、そのポテンシャルエネルギーはまず運動エネルギーに、次に熱と粒子に変換され、宇宙は熱くなる。``再加熱''という名称は、この記述ではやや誤解を招く。なぜなら、それ以前の熱エネルギーはインフレーションの間にすでに希薄化されてしまっているはずだからである。

低い方の極小はゼロにはない。そこに残る小さな値が、現在のゆっくりとした指数関数的膨張に対応するものである。

\begin{classroomqa}
\lecturetimestamp{01:19:48}
\audiencequestion{プラトーは完全に平坦で、摂動が何らかの形で場を端から押し出しているのだろうか。}
\lecturerresponse{いや。完全に平坦ではありえない。何らかの傾きがなければならない。もし平坦すぎれば、再加熱された宇宙は不均一すぎるものになる。}
\end{classroomqa}

``ビッグバン''はここでは近い二つの意味で用いられている。再加熱の出来事そのものを意味することがある。また、最も初期の直接見える段階、すなわち最後の散乱面を意味することもある。確かな点は、実際に観測されるものはすべてインフレーションのプラトーのはるか下にあるということである。

\begin{classroomqa}
\lecturetimestamp{01:45:42}
\audiencequestion{再加熱そのものがビッグバンなのか。}
\lecturerresponse{望むなら、そうだ。しかしその時点で宇宙は依然として極めて高温で不透明であった。最後の散乱面はずっと後、宇宙が十分に冷えて透明になる段階の後に来る。}
\end{classroomqa}

\section{量子揺らぎと銀河の種}

インフレーションは可視宇宙を非常に均一にするが、完全に均一にするわけではない。量子揺らぎがなければ場は完全に滑らかになる。量子力学があれば、場は場所ごとにわずかに変動する。

中心的な図は、通常の空間を縦に、場の値\(\phi\)を横に描いている。もし場が完全に均一であれば、空間の各点は与えられた時刻に同じ\(\phi\)を持ち、各点はそろって端に向かって流れる。するとすべての領域が同じ瞬間に崖から落ちる。量子揺らぎがその完全な同期を乱す。

二つの図が要点を示している。一つは、それ以外は均一な場の中の小さな隆起である。もう一つは、腕を組んで崖に向かって歩く兵士の列である。もし完全に整列していれば、全員が一緒に落ちる。もし何人かがわずかに先に、何人かがわずかに後ろにいれば、異なる時刻に崖を越える。

ここで重要な効果が現れる。より早く落ちた領域はより早く再加熱する。より遅い領域が落ちるころには、より早い領域はすでに膨張して多少冷える時間を持っている。その差が再加熱後のエネルギー密度の変動を生む。

\begin{classroomqa}
\lecturetimestamp{01:34:00}
\audiencequestion{より遅く再加熱された領域は、より早い領域より高温か低温か。}
\lecturerresponse{より早い領域はすでに膨張して冷える時間を持っている。したがってより遅い領域は落ちるときにより高温であり、その差が密度の不均一に寄与する。}
\end{classroomqa}

観測される効果は小さい：
\begin{equation}
\frac{\delta\rho}{\rho}\sim 10^{-5},
\end{equation}
約十万分の一である。これらは原始密度揺らぎであり、暗黒エネルギーではない。この数はここでは定性的なスケッチだけから導かれるものではなく、プラトーがどれほど高く、どれほど平坦でありうるかを制約するために用いられる。

それらはまた必要でもある。それらがなければ銀河は存在しない。再加熱後、重力が初期のパターンを増幅する。過密な領域は過疎な領域から物質を引き寄せる。過密な領域はより過密になり、過疎な領域はより過疎になり、その結果が後の銀河や星の構造である。

これは単なる定性的な話ではない。膨張宇宙における場の量子揺らぎを計算し、それらを再加熱を通して追跡し、その結果を銀河分布から再構成された揺らぎのパターンと比較することができる。一致は非常によい。

同じ議論はエントロピーについての点にも導く。重力による塊形成は直感に反して見えるが、第二法則に反するわけではない。重力は統計力学において特異であり、物事を塊まらせる傾向があり、その塊形成はエントロピーを下げるのではなく上げる。

\begin{classroomqa}
\lecturetimestamp{01:40:24}
\audiencequestion{\(\phi\)は電磁場ではないのに、このすべてにおいてマイクロ波背景はどのように反映されているのか。}
\lecturerresponse{原始密度パターンは最後の散乱面に刻印される。再加熱は場のエネルギーを通常の粒子と光子に変換し、マイクロ波背景は後にそれら原始揺らぎの刻印を運ぶ。}
\end{classroomqa}

\section{音響波と次の段階}

最後に議論は音響波に簡単に移る。これらは後の重力による塊形成より前に来る。原始揺らぎは音響振動に処理され、それらの振動は空に既知の物理的サイズの塊を残す。

それが重要であるのは、既知の距離における既知の物理的サイズが物差しを与えるからである。最後の散乱面でそれが張る角度を測定すれば、幾何学を用いて空間が平坦か、正の曲率か、負の曲率かを知ることができる。詳細は後回しにする。

\begin{classroomqa}
\lecturetimestamp{01:47:19}
\audiencequestion{音響波はどうだろうか。}
\lecturerresponse{それらは天球上の物差しを与える。それらの既知の物理的サイズと、最終散乱面上で観測される角度サイズとを組み合わせることで、空間の幾何学を測定することができるのである。詳細は次回の講義に譲る。}
\end{classroomqa}